\begin{rubric}{Experiência profissional e pesquisa}
\subrubric{\shortstack[l]{Desenvolvedor de Software\\[1pt]{\normalfont\sffamily\small Trixel Desenvolvimento de Sistemas (TRIXEL)}}}
\entry*[2024 -- 2026]%
  \textbf{Projeto VibraTorre} 
  \par Desenvolvimento de soluções de monitoramento da integridade estrutural de torres de transmissão e distribuição, integrando sensores IoT, processamento de sinais, algoritmos de inteligência artificial, simulação estrutural e visualização de dados em tempo real.
  \cvdivider

\subrubric{\shortstack[l]{Desenvolvedor de Visão Computacional\\[1pt]{\normalfont\sffamily\small Instituto Centro de Tecnologia de Software (ICTS)}}}
\entry*[2023 -- 2024]%
  \textbf{Instalação autônoma de esferas sinalizadoras com drones (Projeto ESFERAS)} 
  \par Desenvolvimento de algoritmos de visão computacional e processamento de dados LiDAR para detecção e localização de cabos de linhas de transmissão. Simulação e integração de sistemas robóticos utilizando ROS e Gazebo, com implementação de soluções na plataforma NVIDIA Jetson Xavier NX.
\entry*[2024]%
  \textbf{Reconhecimento facial em hardware de baixo custo (PoC)} 
  \par Desenvolvimento de uma prova de conceito (PoC) de sistema embarcado de reconhecimento facial, utilizando uma imagem Linux personalizada para execução em hardware de baixo custo. Integração de algoritmos de detecção, identificação e teste de vivacidade, com sistema de gerenciamento e cadastro de usuários.
\entry*[2023]%
  \textbf{Contagem de pessoas em ônibus} 
  \par Desenvolvimento de sistema embarcado de visão computacional para detecção e contagem de passageiros, utilizando câmeras e a plataforma NVIDIA Jetson Nano para monitoramento de entradas e saídas em ônibus.
  \cvdivider

\subrubric{\shortstack[l]{Pesquisador PIBIC\\[1pt]{\normalfont\sffamily\small Universidade Federal de Campina Grande (UFCG)}}}
\entry*[2024 -- 2025]%
  \textbf{Aplicação de Redes Neurais Convolucionais Multimodais para Auxiliar a Identificação de Distúrbios Cardíacos através de Dados Estruturados e Não Estruturados.} 
  \par XXII Congresso de Iniciação Científica da UFCG.
\entry*[2022 -- 2023]%
  \textbf{Geração de Imagens Artificiais para Aumento de Dados Utilizando Redes Adversárias Generativas.} 
  \par XX Congresso de Iniciação Científica da UFCG.
\entry*[2021 -- 2022]%
  \textbf{Interpretabilidade e Reusabilidade de Algoritmos Inteligentes para o Diagnóstico Automático de Distúrbios Pulmonares utilizando Radiografias de Tórax e Redes Convolucionais.} 
  \par XIX Congresso de Iniciação Científica da UFCG.
\entry*[2020 -- 2021]%
  \textbf{Classificação de Distúrbios Pulmonares em Radiografias de Tórax Usando Redes Convolucionais.} 
  \par XVIII Congresso de Iniciação Científica da UFCG.
\end{rubric}
